\documentclass[12pt, a4paper]{article}

% --- PACOTES ---
\usepackage[utf8]{inputenc}
\usepackage[portuguese]{babel}
\usepackage{amsmath}               % Pacote principal para matemática
\usepackage{amssymb}               % Símbolos como \mathbb{R}
\usepackage{amsfonts}              % Fontes matemáticas
\usepackage{amsthm}                % Para ambientes de teorema
\usepackage[a4paper, margin=1in]{geometry} % Margens
\usepackage{parskip}               % Espaço entre parágrafos
\usepackage{mathtools}             % Para \coloneqq

% --- CONFIGURAÇÃO DE AMBIENTES ---
\newtheorem{teo}{Teorema}
\newtheorem{prop}{Propriedade}
\newtheorem{defi}{Definição}
\newtheorem{cor}{Corolário}
\newtheorem{macete}{Macete Computacional}

% --- TÍTULO ---
\title{Resumo Acadêmico de Tópicos \\ CE308 – Teoria da Probabilidade 2}
\author{Baseado na Lista 2 do Prof. Benito Olivares Aguilera}
\date{Outubro de 2025}

% --- DOCUMENTO ---
\begin{document}

\maketitle

Este resumo expande os conceitos de \textbf{Estatísticas de Ordem} e \textbf{Distribuição Condicional}, focando nas propriedades formais e estratégias de solução relevantes para a lista.

% ---------------------------------
% SEÇÃO 1: ESTATÍSTICAS DE ORDEM
% ---------------------------------
\section{Distribuição das Estatísticas de Ordem}

Sejam $X_1, X_2, \dots, X_n$ uma amostra aleatória i.i.d. (independentes e identicamente distribuídas) de uma população com f.d.p. (função de densidade de probabilidade) $f_X(x)$ e F.d.a. (função de distribuição acumulada) $F_X(x)$.

As estatísticas de ordem são as variáveis aleatórias ordenadas:
$$ Y_1 \le Y_2 \le \dots \le Y_n $$
Onde $Y_1 = \min(X_1, \dots, X_n)$ e $Y_n = \max(X_1, \dots, X_n)$.

\subsection{Distribuições Marginais (Mínimo, Máximo e k-ésima)}

\begin{macete}
Para o \textbf{Mínimo ($Y_1$)}, use o evento complementar $P(Y_1 > y)$.
Para o \textbf{Máximo ($Y_n$)}, use a definição direta $P(Y_n \le y)$.
\end{macete}

\begin{itemize}
    \item \textbf{Mínimo ($Y_1$):} (Relevante para Q1, Q2, Q4)
        \begin{align*}
            F_{Y_1}(y) &= P(Y_1 \le y) = 1 - P(Y_1 > y) \\
                       &= 1 - P(X_1 > y, \dots, X_n > y) \\
            \text{(i.i.d.)} \quad &= 1 - \prod_{i=1}^n P(X_i > y) = \mathbf{1 - [1 - F_X(y)]^n}
        \end{align*}
        A f.d.p. é obtida pela derivação:
        $$ f_{Y_1}(y) = \frac{d}{dy} F_{Y_1}(y) = \mathbf{n [1 - F_X(y)]^{n-1} f_X(y)} $$

    \item \textbf{Máximo ($Y_n$):} (Relevante para Q1, Q3)
        \begin{align*}
            F_{Y_n}(y) &= P(Y_n \le y) = P(X_1 \le y, \dots, X_n \le y) \\
            \text{(i.i.d.)} \quad &= \prod_{i=1}^n P(X_i \le y) = \mathbf{[F_X(y)]^n}
        \end{align*}
        A f.d.p. é obtida pela derivação:
        $$ f_{Y_n}(y) = \frac{d}{dy} F_{Y_n}(y) = \mathbf{n [F_X(y)]^{n-1} f_X(y)} $$
    
    \item \textbf{k-ésima Estatística ($Y_k$):} (Teoria geral)
        A f.d.p. geral é derivada usando um argumento binomial (quantos $X_i$ são $\le y$):
        $$ f_{Y_k}(y) = \frac{n!}{(k-1)!(n-k)!} [F_X(y)]^{k-1} [1 - F_X(y)]^{n-k} f_X(y) $$
\end{itemize}

\subsection{Distribuições Conjuntas}

\begin{itemize}
    \item \textbf{Conjunta de $Y_1$ e $Y_n$:} (Relevante para Q3)
        Para $y_1 < y_n$, a f.d.p. conjunta é:
        $$ f_{Y_1, Y_n}(y_1, y_n) = \mathbf{n(n-1) [F_X(y_n) - F_X(y_1)]^{n-2} f_X(y_1) f_X(y_n)} $$
        (O termo $[F_X(y_n) - F_X(y_1)]$ é a probabilidade de uma v.a. cair *entre* $y_1$ e $y_n$).

    \item \textbf{Conjunta de todas as $Y_1, \dots, Y_n$:}
        Para $y_1 < y_2 < \dots < y_n$:
        $$ f_{Y_1, \dots, Y_n}(y_1, \dots, y_n) = \mathbf{n! f_X(y_1) f_X(y_2) \cdots f_X(y_n)} $$
        (O $n!$ vem das $n!$ permutações possíveis da amostra original $X_1, \dots, X_n$).
\end{itemize}

\subsection{Casos Notáveis e Propriedades (Macetes)}

\begin{prop}[Caso Exponencial - Q2, Q4]
Se $X_i \sim \text{Exp}(\lambda_i)$ (não necessariamente i.i.d.), então o mínimo (tempo da primeira falha) também é exponencial:
$$ Y_1 = \min(X_1, \dots, X_n) \sim \text{Exp}\left(\sum_{i=1}^n \lambda_i\right) $$
Isto é fundamental em análise de sobrevivência para "riscos concorrentes". Se as $X_i$ são i.i.d. $\text{Exp}(\lambda)$, então $Y_1 \sim \text{Exp}(n\lambda)$.
\end{prop}

\begin{prop}[Caso Uniforme - Q3]
Se $X_i \sim U(0,1)$, então $F_X(x) = x$ e $f_X(x) = 1$ para $x \in (0,1)$. As fórmulas se simplificam drasticamente.
\begin{itemize}
    \item $f_{Y_k}(y) = \frac{n!}{(k-1)!(n-k)!} y^{k-1} (1-y)^{n-k} \sim \mathbf{\text{Beta}(k, n-k+1)}$.
    \item $f_{Y_1, Y_n}(y_1, y_n) = n(n-1)(y_n - y_1)^{n-2}$ para $0 < y_1 < y_n < 1$.
\end{itemize}
Para $U(a,b)$, use a transformação linear $X = a + (b-a)U$ onde $U \sim U(0,1)$.
\end{prop}

\begin{macete}[Cálculo de Covariância - Q3c]
Para calcular $Cov(Y_1, Y_n)$, você precisa de $E[Y_1 Y_n]$, $E[Y_1]$ e $E[Y_n]$.
$$ E[Y_1 Y_n] = \int_a^b \int_{y_1}^b y_1 y_n f_{Y_1, Y_n}(y_1, y_n) \,dy_n \,dy_1 $$
$$ E[Y_k] = \int_a^b y \cdot f_{Y_k}(y) \,dy $$
Usar a distribuição Beta para o caso $U(0,1)$ é muito mais rápido, pois $E[Y_k] = \frac{k}{n+1}$.
\end{macete}

% ---------------------------------
% SEÇÃO 2: CONDICIONAL
% ---------------------------------
\section{Distribuição e Esperança Condicional}

Este tópico trata de como a informação sobre uma v.a. $X$ altera nosso conhecimento sobre $Y$.

\subsection{Definições Formais}

\begin{defi}[Distribuição Condicional]
Sejam $(X,Y)$ v.a. contínuas com densidade conjunta $f(x,y)$ e marginais $f_X(x)$ e $f_Y(y)$.
\begin{itemize}
    \item \textbf{Marginal de X:} $f_X(x) = \int_{-\infty}^{\infty} f(x,y) \,dy$
    \item \textbf{Condicional de Y dado X:} $f(y|x) = \dfrac{f(x,y)}{f_X(x)}$, definida para $f_X(x) > 0$.
\end{itemize}
(Análogo para o caso discreto $p(y|x) = p(x,y)/p_X(x)$ - Q10, Q13, Q14).
\end{defi}

\begin{macete}[Modelo Hierárquico - Q9, Q12, Q18, Q20]
Muitos problemas (especialmente Bayesianos) são definidos em etapas:
\begin{enumerate}
    \item $X$ é sorteado de $f_X(x)$ (distribuição "a priori").
    \item $Y$ é sorteado de $f(y|x)$ (distribuição "verossimilhança").
\end{enumerate}
Para resolver, use as "leis" na ordem inversa:
\begin{itemize}
    \item \textbf{Achar Conjunta:} $f(x,y) = f(y|x) f_X(x)$
    \item \textbf{Achar Marginal de Y:} $f_Y(y) = \int_{-\infty}^{\infty} f(x,y) \,dx = \int f(y|x) f_X(x) \,dx$
    \item \textbf{Achar $f(x|y)$ (Regra de Bayes):} $f(x|y) = \dfrac{f(x,y)}{f_Y(y)} = \dfrac{f(y|x) f_X(x)}{\int f(y|u) f_X(u) \,du}$
\end{itemize}
\end{macete}

\subsection{Esperança Condicional}

\begin{defi}[Esperança Condicional]
\begin{enumerate}
    \item $E[Y|X=x]$ é um \textbf{número} (ou função de $x$), calculado como a esperança de $Y$ usando a densidade $f(y|x)$:
    $$ E[Y|X=x] \coloneqq \int_{-\infty}^{\infty} y \cdot f(y|x) \,dy $$
    \item $E[Y|X]$ é uma \textbf{Variável Aleatória}. É a função $g(x) = E[Y|X=x]$ aplicada à v.a. $X$, ou seja, $g(X)$.
\end{enumerate}
Formalmente, $E[Y|X]$ é a única v.a. $\sigma(X)$-mensurável que satisfaz $E[E[Y|X] \cdot \mathbb{I}(X \in A)] = E[Y \cdot \mathbb{I}(X \in A)]$ para todo conjunto $A$.
\end{defi}

\subsection{Propriedades Fundamentais da Esperança Condicional}

\begin{prop}[Linearidade]
$E[aY + bZ | X] = aE[Y|X] + bE[Z|X]$
\end{prop}

\begin{prop}[Tirar o que é Conhecido (Taking out what is known)]
Se $g$ é uma função, $g(X)$ é "conhecido" dado $X$:
$$ E[g(X) \cdot Y | X] = g(X) \cdot E[Y|X] $$
\end{prop}

\begin{prop}[Independência - Q11]
Se $X$ e $Y$ são independentes ($X \perp Y$), então $E[Y|X] = E[Y]$ (constante). A informação sobre $X$ não altera a esperança de $Y$.
\end{prop}

\begin{prop}[Idempotência (Lei da Torre)]
$E[ E[Y|X] | X ] = E[Y|X]$ (A melhor estimativa da melhor estimativa é a própria melhor estimativa).
\end{prop}

\begin{prop}[Projeção $L^2$ - Q22]
$E[Y|X]$ é a "melhor" aproximação de $Y$ que pode ser feita usando apenas $X$, no sentido de que minimiza a variância do erro (Mínimos Quadrados):
$$ E[Y|X] = \arg \min_{g(X)} E[(Y - g(X))^2] $$
O erro $W = Y - E[Y|X]$ é \textbf{ortogonal} a qualquer função de $X$:
$$ E[W \cdot g(X)] = 0 \implies \mathbf{Cov(X, Y - E[Y|X]) = 0} $$
(Isto prova a Q22).
\end{prop}

\subsection{Teoremas de Decomposição (As Leis Totais)}

\begin{teo}[Lei da Esperança Total - Q19, Q20, Q21, Q24]
A esperança da esperança condicional é a esperança original.
$$ \mathbf{E[Y] = E[E[Y|X]]} $$
\end{teo}

\begin{teo}[Lei da Variância Total - Q19, Q20, Q23]
A variância total é a soma da esperança da variância condicional e a variância da esperança condicional (Eve's Law: EVVE).
$$ \mathbf{Var(Y) = E[Var(Y|X)] + Var[E[Y|X]]} $$
\end{teo}

\begin{teo}[Lei da Covariância Total - Q27]
$$ \mathbf{Cov(X,Y) = E[Cov(X,Y|Z)] + Cov[E(X|Z), E(Y|Z)]} $$
\end{teo}

\begin{macete}[Cálculo de Variância via Leis Totais - Q20]
Para calcular $Var(Y)$ usando a Lei Total (LTV):
\begin{enumerate}
    \item Calcule $g(X) = E[Y|X]$.
    \item Calcule $h(X) = Var(Y|X)$.
    \item Calcule $Var(Y) = E[h(X)] + Var[g(X)]$.
\end{enumerate}
Frequentemente, é mais fácil calcular $Var(Y)$ usando $E[Y^2] - (E[Y])^2$:
\begin{itemize}
    \item $E[Y] = E[E[Y|X]]$ (usando a LTE).
    \item $E[Y^2] = E[E[Y^2|X]]$ (usando a LTE na v.a. $Y^2$).
    \item \textbf{Truque:} $E[Y^2|X] = Var(Y|X) + (E[Y|X])^2 = h(X) + [g(X)]^2$.
    \item Combine tudo: $Var(Y) = E[h(X) + g(X)^2] - (E[g(X)])^2$.
\end{itemize}
\end{macete}

\subsection{Modelos Conjugados e Casos Notáveis da Lista}

\begin{itemize}
    \item \textbf{Poisson-Binomial (Q12, Q13):} Se $Y \sim \text{Poisson}(\lambda)$ (nº de passas total) e $X|Y=n \sim \text{Bin}(n, p)$ (nº de passas que você recebe), então a marginal de X é:
    $$ X \sim \text{Poisson}(\lambda p) $$
    (Intuitivo: "afinamento" (thinning) de um processo de Poisson).

    \item \textbf{Binomial-Binomial (Q14):} Se $X \sim \text{Bin}(m,p)$, $Y \sim \text{Bin}(n,p)$, $S = X+Y$.
    Então $S \sim \text{Bin}(m+n, p)$. A condicional é:
    $$ X | S=k \sim \text{Hipergeométrica}(N=m+n, K=m, n=k) $$
    Portanto, $E[X|S=k] = k \cdot \frac{m}{m+n} \implies E[X|S] = S \cdot \frac{m}{m+n}$.

    \item \textbf{Geométrica-Geométrica (Q16):} Se $X_1, X_2 \sim \text{Geom}(p)$ (i.i.d.), $S = X_1+X_2$.
    Então $S \sim \text{Binomial Negativa}(r=2, p)$. A condicional é:
    $$ X_1 | S=k \sim \text{Uniforme Discreta}\{0, 1, \dots, k\} $$
    (A probabilidade é $1/(k+1)$ para cada ponto).
\end{itemize}

\end{document}